\documentclass[10pt,a4paper]{amsart}
\usepackage[margin=19mm]{geometry}
\usepackage{amsmath,amssymb,mathtools}
\usepackage[scr=rsfs,cal=euler]{mathalfa}
\usepackage{xfrac}
\usepackage[unicode,hidelinks,bookmarks=false]{hyperref}
\newcommand{\ie}{i.e.\ }
\newcommand{\cf}{cf.\ }
\newcommand{\resp}{respectively\ }
\newcommand{\Subsection}{\subsection}
\providecommand{\nobreakdash}{\nobreak\hbox{-}}
\setlength{\emergencystretch}{2em}
\multlinegap0pt
\usepackage{bm}
\usepackage{bbm}
\usepackage{dsfont}
\usepackage{xspace}
\usepackage{cancel}
\usepackage{mathtools}
\usepackage{amsthm}
\makeatletter
\@ifundefined{theorem}{\newtheorem{theorem}{Theorem}}{}
\makeatother
\makeatletter
\newcommand{\NN}{\mathbb{N}}
\newcommand{\PP}{\mathds{P}}
\newcommand{\RR}{\mathbb{R}}
\expandafter\ifx\csname assumption\endcsname\relax
\newenvironment{assumption}[2]
{%
	\renewcommand{\assumptionnumber}{(\textbf{#1#2})}%
	\begin{assumption*}%
		\protected@edef\@currentlabel{(\textbf{#1#2})}%
	}
\fi
\newcommand{\wt}{\widetilde}
\makeatother
\title[Modelling Anomalous Diffusion: The Role of CTRWs and Non-Loc]{Modelling Anomalous Diffusion: The Role of CTRWs and Non-Local Dynamics}
\author{Ivan Bio\v{c}i\'{c}, Bruno Toaldo}
\date{Source: arXiv:2607.27150v1 (CC BY 4.0)}
\begin{document}
\maketitle

\noindent\emph{Source and licence.} Modelling Anomalous Diffusion: The Role of CTRWs and Non-Local Dynamics. Version arXiv:2607.27150v1 (CC BY 4.0), distributed under the Creative Commons Attribution 4.0 International licence, as recorded in the arXiv licence metadata for this exact version and shown on the abstract page https://arxiv.org/abs/2607.27150v1. Full rights notice: licenses/arxiv-2607.27150-CC-BY-4.0.txt.\par\smallskip

\noindent\textit{Editorial scope.} Counted unit: the Theorem whose source label is \texttt{thmGS+1} in the article (source0.tex lines 405--412, source label \texttt{thmGS+1}), delivered together with the article's own complete original proof of it (source0.tex lines 413--439, one \texttt{proof} environment of 27 lines). No mathematics is written, repaired or re-derived: statement and proof are the article's own text line for line. Required context retained uncounted: defctrw (lines 206--209); nonexplosion (lines 222--225); thmGS (lines 233--257); thmGS+ (lines 356--372); next-jump (lines 359--361). Every definition, display and label of those lines is retained unchanged. Omitted from this package and not redistributed: 1--205, 210--221, 226--232, 258--355, 362--404, 440--2318. Nothing inside a retained excerpt is omitted or altered: every retained body is delivered exactly as the article's lines are written, so no omission receipt of any kind accompanies this package. Citations: the retained excerpts carry no citation command at all, so no bibliography entry of the article is transcribed and no reference is redistributed. Reference adaptation: none was needed and none was made; every reference inside the delivered unit resolves inside it, so no unresolved reference or citation is produced. Printed slips kept literal: no printed word, symbol, hypothesis or number of the counted statement or of its proof was altered. Layout and class: the article's own class is \texttt{amsart} with packages including amssymb, amsmath, amsfonts, natbib, lscape, upgreek, aligned-overset, color, comment, mathtools, dsfont, mathrsfs, bm, graphicx; the wrapper is the lane's pinned \texttt{amsart} editorial wrapper at 10pt on A4, which restates the theorem-like environments the retained text needs and adds the article's own macro definitions that the retained text needs (\texttt{\textbackslash NN}, \texttt{\textbackslash PP}, \texttt{\textbackslash RR}, \texttt{\textbackslash wt}), with extra wrapper packages (bm, bbm, dsfont, xspace, cancel, mathtools). The delivered numbering is the unit's own. Assets: no retained span contains a figure environment, a graphics-inclusion command, a PDF-page inclusion, a table or a TikZ picture; no image file is copied into this package, the package declares no asset dependency and no image file is redistributed with it. Licence: version arXiv:2607.27150v1 of this article is distributed under the Creative Commons Attribution 4.0 International licence, as recorded in the arXiv licence metadata for this exact version and shown on the abstract page https://arxiv.org/abs/2607.27150v1. A full rights notice is delivered at licenses/arxiv-2607.27150-CC-BY-4.0.txt. Characters: every retained excerpt is pure ASCII, so the delivered engine, XeLaTeX with its default Latin Modern text font, reports no missing character.
	\begin{align}
		Y(t) \coloneqq S_{N(t)}.
		\label{defctrw}
	\end{align}

	\begin{align}
		\sum_n W_n = +\infty, \quad \PP^{(x,s)}\text{--a.s.},
		\label{nonexplosion}
	\end{align}

	\begin{theorem}
		\label{thmGS}
		Let $Y(t)$, $t\in\RR$, be a CTRW as in \eqref{defctrw} and suppose that $(S_n, T_n)$, $n\in \NN \cup \{0\}$, is a Markov chain such that \eqref{nonexplosion} holds. Assume further that, for all $(z,r) \in \mathbb{R}^d \times \mathbb{R}$, $(z^\prime,r^\prime) \in \mathbb{R}^d \times \mathbb{R}$ and $(x,s) \in \mathbb{R}^d \times \mathbb{R}$, it holds that
		\begin{align}\label{semi-jump}
			\mathds{P}^{(z,r)} \left( S_{n+1} \in dy, W_{n+1} \in dw \mid S_n = x, T_n = s \right) \, = \, &\mathds{P}^{(z^\prime,r^\prime)} \left( S_{n+1} \in dy, W_{n+1} \in dw \mid S_n = x \right) \notag \\
			= \, & \mathds{P}^{(x,0)} \left( S_1 \in dy, W_1 \in dw \right) \notag  \\
			\eqqcolon \, & \mu_x(dy,dw).
		\end{align}
		Define, for $t \geq 0$,
		\begin{align}
			\xi(t) \, = \, t-T_n, \qquad T_n\leq t < T_{n+1},
		\end{align}
		i.e. $\xi(t) = t- T(t)$.
		Then the process $(Y(t), \xi (t))$, $t \geq 0$, is a time-homogeneous Markov process, under the measures $\PP^{(x,s)}$, $(x,s)\in \RR^d\times (-\infty,0]$, and the filtration $\mathcal{F}_t$, $t \geq 0$, generated by $N(s), S_0, \dots, S_{N(s)}, T_0, \dots, T_{N(s)}$, $s \le t$. Moreover, its semigroup $P_t$, $t \geq 0$, satisfies for $(x,s)\in \RR^d\times [0,+\infty)$
		\begin{align}
			P_tf(x,s) \, = \, f(x,s+t) \, \frac{\mu_x(\mathbb{R}^d, (t+s, +\infty))}{\mu_x(\mathbb{R}^d, (s, +\infty))} + \int_{(s,s+t]} \int_{\mathbb{R}^d} P_{t+s-u}f(y,0)   \frac{\mu_x(dy, du)}{\mu_x(\mathbb{R}^d, (s, +\infty))},
			\label{transition_step}
		\end{align}
		and 
		\begin{align}
			P_tf(x,0) = \mathds{E}^{(x,0)} f(Y(t), \xi(t)), \quad x\in \RR^d,
			\label{semizero}
		\end{align}
		for any function $f$ which is bounded and Borel measurable ($f \in \mathcal{B}_b (\mathbb{R}^d \times [0, +\infty))$).
	\end{theorem}

	\begin{theorem}
		\label{thmGS+}
		Suppose that the assumptions of Theorem \ref{thmGS} are satisfied. Define
		\begin{align}\label{next-jump}
			\Xi(t) \, = \, T_{n+1}-t, \qquad T_n \le  t < T_{n+1},
		\end{align}
		i.e. $\Xi (t) = T^+(t)-t$, $t \geq 0$. Denote by $N^-(s):= N(s-)$, $s \geq 0$.
		Then the process $(Y^+(t-), \Xi(t-))$, $t \geq 0$, is a time-homogeneous Markov process, under all $\PP^{(x,s)}$, $(x,s)\in \RR^d\times (-\infty,0]$, and the filtration $\mathcal{G}^-_t$, $t \geq 0$, generated by $N^-(s)$, $S_0, \dots, S_{N^-(s)+1}$, $T_0, \dots, T_{N^-(s)+1}$, $s \leq t$. Its semigroup $P_t^+$, $t \geq 0$, satisfies for $t\ge 0$ and $(x,s)\in\RR^d\times [0,+\infty)$
		\begin{align}\label{710-a1}
			P_t^+f(x,s) \, = \, \mathds{1}_{[0 \leq t \le  s]} f(x,s-t)+ \mathds{1}_{[0 \leq s <t]} P_{t-s}^+f(x,0),
		\end{align}
		and
		\begin{align}\label{710-b1}
			P_t^+f(x,0) \, = \, \mathds{E}^{(x,0)}f(Y(t-), \Xi(t-)),\quad t>0,\, x\in \RR^d.
		\end{align}
		for any function $f \in \mathcal{B}_b (\mathbb{R}^d \times [0, +\infty))$.
	\end{theorem}

	\begin{theorem}
		\label{thmGS+1}
		Suppose that the assumptions of Theorem \ref{thmGS} are satisfied, and let $\Xi(t)$, $t\in \RR$, be defined as in \eqref{next-jump}. Then the process $(Y^+(t), \Xi(t))$, $t \geq 0$, is a time-homogeneous Markov process, under all $\PP^{(x,s)}$, $(x,s)\in \RR^d\times (-\infty,0]$, and the filtration $\mathcal{G}_t$, $t \geq0$, generated by $N(s), S_0, \dots, {S_{N(s)+1}}, T_0, \dots, {T_{N(s)+1}}$, $s \leq t$. Its semigroup $\wt P_t^+$, $t \geq 0$, satisfies for $(x,s)\in \RR^d\times (0,+\infty)$
		\begin{align}\label{710-a}
			\wt P^+_tf(x,s)&=\mathds{1}_{[0\le t<s]}f(x,s-t)+\mathds{1}_{[0<s\le t ]}\int_{\RR^d}\int_{(0,+\infty)}\wt P^+_{t-s}f(y,w)\mu_{x}(dy,dw),
		\end{align}
		for any function $f \in \mathcal{B}_b (\mathbb{R}^d \times (0, +\infty))$.
	\end{theorem}

	\begin{proof}
		The proof is again very similar to the proofs of Theorem \ref{thmGS} and Theorem \ref{thmGS+}. However, since the claim also appears to be new, we give a detailed sketch.
		
		We have that
		\begin{align}
			&\mathds{E}^{(x,s)} \left[ f(Y^+(T+t), \Xi(T+t)) \mid \mathcal{G}_T \right] \notag \\
			\begin{split}\label{1612}
				\, = \, & \mathds{E}^{(x,s)} \left[ f(Y^+(T+t), \Xi(T+t)) \mathds{1}_{[T+t < T_{N(T)+1}]} \mid \mathcal{G}_T \right] \notag \\ & + \mathds{E}^{(x,s)} \left[ f(Y^+(T+t), \Xi(T+t)) \mathds{1}_{[T+t \geq T_{N(T)+1}]} \mid \mathcal{G}_T \right].
			\end{split}
		\end{align}
		The first term in \eqref{1612} is just
		\begin{align}
			\mathds{E}^{(x,s)} \left[ f(Y^+(T+t), \Xi(T+t)) \mathds{1}_{[T+t < T_{N(T)+1}]} \mid \mathcal{G}_T \right] \, = \, f(Y^+(T), \Xi(T)-t)\mathds{1}_{[t < \Xi(T)]}.
		\end{align}
		For the second term, instead, we obtain
		\begin{align}
			&\mathds{E}^{(x,s)} \left[ f(Y^+(T+t), \Xi(T+t)) \mathds{1}_{[T+t \geq T_{N(T)+1}]} \mid \mathcal{G}_T \right] \notag \\
			= \, &\mathds{E}^{(x,s)} \left[ \mathds{1}_{[T+t \geq T_{N(T)+1}]} \mathds{E}^{(x,s)} \left[  f(Y^+(T+t), \Xi(T+t)) \mid \mathcal{G}_T, T_{N(T)+2}, S_{N(T)+2} \right] \mid \mathcal{G}_T  \right].\label{1142}
		\end{align}
		Now observe that the process $(Y^+(T+t), \Xi(T+t))$, on $T+t \geq T_{N(T)+1}$, conditionally on $\mathcal{G}_T$, $S_{N(T)+2}$, $T_{N(T)+2}$ is constructed with respect to the r.v.'s $S_{N(T)+3}$, $T_{N(T)+3},\dots$, in the same way the process $(Y^+(T+t-T_{N(T)+1}), \Xi(T+t-T_{N(T)+1}))$ is constructed with respect to the r.v.'s $S_2$, $T_2$, $\dots$, conditionally on $T_0=0$, $S_1 = S_{N(T)+2}$, and $T_1 = T_{N(T)+2}-T_{N(T)+1}$. Therefore \eqref{1142} becomes
		\begin{align}
			&\mathds{E}^{(x,s)} \left[ \mathds{1}_{[T+t \geq T_{N(T)+1}]} \mathds{E}^{(x,0)} \left[  f(Y^+(T+t-T_{N(T)+1}), \Xi(T+t-T_{N(T)+1})) \right. \right. \notag  \\ & \left. \left. \qquad \mid S_1=S_{N(T)+2}, T_1= T_{N(T)+2}-T_{N(T)+1} \right] \mid \mathcal{G}_T  \right] \notag \\
			= \, &\mathds{1}_{[T+t \geq T_{N(T)+1}]} \int_0^{+\infty} \int_{\mathbb{R}^d} \wt P^+_{T+t-T_{N(T)+1}}f(y, w) \mathds{P}^{(x,s)} (S_{N(T)+2} \in dy, W_{N(T)+2} \in dw \mid \mathcal{G}_T) \notag \\
			= \, &\mathds{1}_{[t \geq \Xi(T)]} \int_0^{+\infty} \int_{\mathbb{R}^d} {\wt P}^+_{t-\Xi(T)}f(y, w) \mu_{{Y^+(T)}} (dy, dw).
		\end{align} 
		This concludes the proof.
	\end{proof}

\end{document}
